\begin{frame}
Alcuni strumenti analizzati sono: 
\begin{itemize} 
\item Virtualbox 
\item Scapy 
\item RITA 
\end{itemize}
Invece alcuni progetti analizzati sono: 
\begin{itemize} 
\item ICMP door 
\item ICMP exfil 
\item ICMP tunnel 
\end{itemize}
\end{frame}
%
\begin{frame}
\frametitle{Strumenti utilizzati}
VirtualBox è un software di virtualizzazione open source che consente di eseguire più sistemi operativi (guest) su un singolo computer (host). 
\MakeUppercase{è} stato usato per la creazione di macchine virtuali Ubuntu; così da testare il codice sviluppato in quell'ambiente. 
%Per ciascun entità sarà quindi presente una macchina virtuale. 
%Alla fine si sono ottenute quattro macchine virtuali: una per l'attaccante e la vittima, mentre le altre due per i proxy. 
%Si poteva usare anche un singolo proxy ma si voleva testare come i dati ricavati dalla vittima (e da inoltrare all'attaccante), venissero distribuiti ai proxy disponibili.  
\vspace{1ex} \newline 
Scapy è un framework per la manipolazione dei pacchetti scritto in Python. 
Consente di falsificare varie tipologie di pacchetti %(http, tcp, ip, udp, icmp, ecc.). 
%ed è in grado di creare o decodificare pacchetti di vari protocolli. 
Inoltre può inviarli in rete, catturarli, memorizzarli o leggerne i dati. 
La presenza di metodi che permettono l'inio e la fine dello sniffing; ha permesso lo sviluppo dei canali di comunicaizone tramite protocollo ICMP. 
%Nella libreria il meotdo \textbf{send} (o similari) permetteranno di inviare un definito pacchetto. 
%Invcece per definire un pacchetto si concatenano i livelli \cite{scapy-inet} indicanti i protocolli presenti (nel nostro caso IP e ICMP). per poter ascoltare il traffico di rete \cite{scapy-sniffing} è presente il metodo \textbf{sniff} oppure si può definire una variabile di tipo \textit{AsyncSniffer}; 
\vspace{1ex} \newline  
RITA è uno strumento open source per la ricerca delle 
minacce di rete, progettato per identificare attività 
di comando e controllo (C2) dannose. Acquisisce i log 
di Zeek e utilizza l'analisi comportamentale per 
identificare sistemi potenzialmente compromessi. 
\vspace{1ex} \newline 
Una costante assoluta su cui RITA fà affidamento per rilevare i malware, è che dovranno chiamare "casa". 
Da questo presupposto; analizzando il traffico di rete, rilevare le chiamate C2 indipendentemente dalla piattaforma. 
La persistenza è l'attributo chiave che si ricerca quando si analizza sistemi compromessi. 
RITA cerca gli indicatori principali relativi a questa persistenza. 
Viene fatto acquisendo i log di connessione di Zeek e tramite l'utilizzo dell'analisi comportamentale 
identifica i sistemi potenzialmente compromessi. 
\vspace{2ex} \newline 
Le funzionalità principali sono: 
il rilevamento dei Beacon, il rilevamento del tunneling DNS, il rilevameno di connessioni che hanno comunicato per tempi lunghi, 
controllo dei feed per le Threat Intel (domini e host sospetti), la valutazione per gravità delle connessioni, quanti host hanno comunicato 
con un determinato host, il primo incontro di un host. 
\end{frame} 
%
\begin{frame}
\frametitle{ICMP Door}
\MakeUppercase{è} stato studiato per comprendere come 
poter strutturare il tunneling dei dati. 
Oltre alla struttura delle entità, che successivamente 
verranno ridefinite, risulta interessante che il 
programma richieda degli argomenti dall'utente. 
Tramite la libreria \textit{argparse} richiede 
l'interfaccia su cui ascoltare i dati e l'indirizzo 
di destinazione dei pacchetti. 
Inoltre tramite il metodo \textit{sniff} del framework 
Scapy ascolta il flusso dei dati mentre  con \textit{sr} 
invia i paccheti. 
\vspace{2ex} \newline 
%Il programma introduce a una possibile struttura del 
%Covert Channel; 
Alcuni svantaggi saranno che i dati vengono trasmessi in chiaro 
nel campo Datra del messaggio ICMP Echo Reply. Inoltre 
il campo identifier rimmarrà invariato per tutta la 
comunicazione e siccome a molteplici richieste non 
potrà essere associata una richiesta. Un sistema di sicurezza 
potrebbe bloccare i pacchetti. 
\end{frame}
%
\begin{frame}
\frametitle{ICMP Exfil}
Si è analizzato ICMP Exfil per vedere come un Covert Channel temporalizzato postesse funzionare. 
Riceve un dato, lo converte in binario e dopodichè avrà una lista di numeri binari. 
Per mandare il dato, invia un ping con un timeout pari a \textbf{binary\_number+leway}. 
Il valore leway viene utilizzato per rallentare il numero di pacchetti inviati ed avere una connesisone maggiormente silenziosa. 
L'autore infine indica come miglioria, la crittografia dei dati per aggiungere del rumore, dell'entropia. 
\vspace{1ex} \newline 
Ciò su cui si affida è che un osservatore, vedendo i pacchetti ICMP, li veda come vailidi; 
Siccome non riuscirebbero a trovare alcun dato, a meno che non sappiano della tecnica utilizzata. 
\end{frame}
%
\begin{frame}
\frametitle{ICMP Tunnel}
\MakeUppercase{è} uno strumento che permette il tunneling del traffico IP. 
Tramite delle richieste e risposte ICMP Echo, incapsula il traffico e lo invia al server proxy. 
Quest'ultimo lo decapsulano e lo inoltra. 
I pacchetti in entrata, che sarebbero diretti alla macchina vittima, sono poi incapsulati dal proxy e inviati. 
L'approccio è possibile siccome RFC-792, che indica le linee guida del protocollo ICMP, permette una quantità arbitraria di dati nei pacchetti ICMP Echo (sia richiesta che risposta). 
\end{frame} 